\documentclass{article}
\usepackage{amsmath}
\usepackage{amssymb}
\usepackage{enumitem}

\begin{document}

\title{Lecture Scribe \\ CSE 400 – Fundamentals of Probability in Computing \\ Tutorial 2 Solutions \\ School of Engineering and Applied Science (SEAS), Ahmedabad University}
\maketitle

This lecture scribe reconstructs the reasoning and derivations used in solving the problems presented in Tutorial 2. Each problem is developed step-by-step so that the logical progression of the arguments can be followed during exam revision.

\section{Q1. Random Point on a Line Segment}

\subsection{Problem Statement}
A point is chosen randomly on a line segment of length $L$. The line is divided into two segments by this point. We must determine the probability that the ratio of the shorter segment to the longer segment is less than $1/4$.

\subsection{Step 1: Mathematical Representation}
Let $x$ be the distance of the chosen point from one endpoint of the line. \\
Since the point lies on the line segment, \\
$0 < x < L$ \\
The two resulting segments are: \\
First segment: $x$ \\
Second segment: $L-x$

\subsection{Step 2: Ratio of Segments}
The ratio of the shorter segment to the longer segment is \\
$R = \frac{\min(x, L-x)}{\max(x, L-x)}$ \\
We want to find the probability that \\
$R < \frac{1}{4}$

\subsection{Step 3: Case Analysis}
The nature of the shorter segment depends on whether the point lies to the left or right of the midpoint $L/2$.

\subsubsection{Case 1: $x \le \frac{L}{2}$}
In this case: \\
shorter segment = $x$ \\
longer segment = $L-x$ \\
Thus the condition becomes \\
$\frac{x}{L-x} < \frac{1}{4}$ \\
Multiplying both sides by $4(L-x)$: \\
$4x < L-x$ \quad $5x < L$ \quad $x < \frac{L}{5}$ \\
Therefore, the valid values in this case are \\
$0 < x < \frac{L}{5}$

\subsubsection{Case 2: $x > \frac{L}{2}$}
Now: \\
shorter segment = $L-x$ \\
longer segment = $x$ \\
Thus \\
$\frac{L-x}{x} < \frac{1}{4}$ \\
Multiplying both sides by $4x$: \\
$4L - 4x < x$ \quad $5x > 4L$ \quad $x > \frac{4L}{5}$ \\
Thus the valid interval becomes \\
$\frac{4L}{5} < x < L$

\subsection{Step 4: Valid Region}
Combining both cases: \\
$(0, L/5) \cup (4L/5, L)$ \\
Length of valid region: \\
$\frac{L}{5} + \frac{L}{5} = \frac{2L}{5}$

\subsection{Step 5: Probability Calculation}
Because the point is uniformly distributed along the segment, \\
$P = \frac{\text{valid region}}{\text{total region}}$ \quad $P = \frac{2L/5}{L}$ \quad $P = \frac{2}{5}$

\section{Q2. Bayesian Decision Problem with Normal Distributions}
Given \\
Two regions exist in an image: \\
\textbf{White Region} \\
$X|W \sim N(4,4)$ \\
Thus \\
$\sigma_1 = 2$

\textbf{Black Region} \\
$X|B \sim N(6,9)$ \\
Thus \\
$\sigma_2 = 3$

\textbf{Prior Probabilities} \\
$P(B) = \alpha$ \quad $P(W) = 1 - \alpha$ \\
A measurement is observed: \\
$X = 5$ \\
The problem asks for the value of $\alpha$ such that the probability of classification error is the same regardless of the decision made.

\subsection{Step 1: Equal Posterior Probabilities}
Equal classification error implies \\
$P(B|X=5) = P(W|X=5) = \frac{1}{2}$

\subsection{Step 2: Apply Bayes' Theorem}
$P(B|X=5)= \frac{\alpha f(5|B)} {\alpha f(5|B)+(1-\alpha)f(5|W)}$ \\
Setting this equal to $1/2$ gives the equation needed to solve for $\alpha$.

\subsection{Step 3: Normal Density Evaluation}
The normal probability density function is \\
$f(x)=\frac{1}{\sigma\sqrt{2\pi}}e^{-\frac{(x-\mu)^2}{2\sigma^2}}$ \\
Substituting parameters: \\
$f(5|B)=\frac{1}{3\sqrt{2\pi}}e^{-1/18}$ \quad $f(5|W)=\frac{1}{2\sqrt{2\pi}}e^{-1/8}$

\subsection{Step 4: Substitution}
Substituting into Bayes' expression: \\
$\frac{\alpha e^{-1/18}}{\alpha e^{-1/18}+(1-\alpha)e^{-1/8}\frac{3}{2}}=\frac{1}{2}$ \\
Simplifying yields \\
$\frac{3\alpha}{3\alpha+2(1-\alpha)e^{-5/72}}=\frac{1}{2}$

\subsection{Step 5: Solve for $\alpha$}
Using the approximation \\
$e^{-5/72} \approx 0.93$ \quad $6\alpha = 3\alpha + 2(1-\alpha)(0.93)$ \quad $3\alpha = 1.86 - 1.86\alpha$ \quad $4.86\alpha = 1.86$ \quad $\alpha \approx 0.3827$ \\
Thus \\
$\alpha \approx 0.38$

\section{Q3. Optimal Placement of a Fire Station}

\subsection{Part (a): Finite Road}
Let \\
$X \sim \text{Uniform}(0,A)$ \\
with density \\
$f_X(x)=\frac{1}{A}$ \\
The objective is to minimize \\
$E[|X-a|]$

\subsubsection{Step 1: Express the Expectation}
$E[|X-a|] = \int_0^a (a-x)\frac{1}{A}dx + \int_a^A (x-a)\frac{1}{A}dx$

\subsubsection{Step 2: Differentiate}
The resulting expression simplifies to \\
$\frac{2a^2 - 2aA + A^2}{2A}$ \\
Differentiating with respect to $a$: \\
$\frac{1}{2A}(4a-2A)$

\subsubsection{Step 3: Minimization}
Set derivative to zero: \\
$4a-2A=0$ \quad $a=\frac{A}{2}$ \\
Therefore, the station should be located at the midpoint of the road.

\subsection{Part (b): Infinite Road}
Now \\
$X \sim \text{Exponential}(\lambda)$ \\
PDF: \\
$f_X(x)=\lambda e^{-\lambda x}$ \\
The expected distance is minimized when \\
$\frac{d}{da}E[|X-a|]=1-2e^{-\lambda a}$ \\
Setting this equal to zero: \\
$e^{-\lambda a}=\frac{1}{2}$ \\
Taking logarithms: \\
$-\lambda a=\ln(1/2)$ \quad $a=\frac{\ln 2}{\lambda}$

\section{Q4. Mixture of Exponential Battery Lifetimes}
Let $X$ denote battery lifetime. \\
Two types exist with probabilities $p_1$ and $p_2$. \\
Goal: \\
$P(X>t+s|X>t)$ \\
Using total probability: \\
$P(X>t)=p_1e^{-\lambda_1 t}+p_2e^{-\lambda_2 t}$ \quad $P(X>t+s)=p_1e^{-\lambda_1(t+s)}+p_2e^{-\lambda_2(t+s)}$ \\
Thus \\
$P(X>t+s|X>t)= \frac{p_1e^{-\lambda_1(t+s)}+p_2e^{-\lambda_2(t+s)}} {p_1e^{-\lambda_1 t}+p_2e^{-\lambda_2 t}}$

\section{Q5. Conditional Distribution of Poisson Variables}
Let \\
$X\sim \text{Poisson}(\lambda_1)$ \quad $Y\sim \text{Poisson}(\lambda_2)$ \\
Independent. \\
We seek \\
$P(X=k|X+Y=n)$ \\
Using conditional probability: \\
$P(X=k|X+Y=n)= \frac{P(X=k,Y=n-k)}{P(X+Y=n)}$ \\
Substituting Poisson PMFs and simplifying results in \\
$P(X=k|X+Y=n)= \binom{n}{k} \left(\frac{\lambda_1}{\lambda_1+\lambda_2}\right)^k \left(\frac{\lambda_2}{\lambda_1+\lambda_2}\right)^{n-k}$ \\
Thus the conditional distribution is binomial.

\section{Remaining Problems (Conceptual Summary)}

\subsection{Q6}
Transformation of a uniform random variable producing a distribution with: \\
continuous density in certain intervals \\
a discrete point mass at $Y=0$ equal to $1/2$.

\subsection{Q7}
Stock price movements are modeled with independent Bernoulli trials, producing a Binomial distribution for the number of upward moves.

\subsection{Q8}
Repair time follows an Exponential distribution with parameter $\lambda = 1/2$. \\
Key results include: \\
$P(X>2)=e^{-1}$ \\
and \\
$P(X\ge10|X>9)=e^{-1/2}$

\subsection{Q9}
For \\
$Z=\sqrt{X^2+Y^2}$ \\
the CDF is computed by integrating over a circular region $x^2+y^2 \le z^2$. \\
The PDF is obtained by differentiating the CDF using Leibniz's rule.

\subsection{Q10}
For the maximum of i.i.d. variables \\
$M_n=\max(X_1,...,X_n)$ \\
the joint distribution of $M_n$ and $M_{n+1}$ depends on whether $x\le y$ or $x>y$: \\
$P(M_n\le x,M_{n+1}\le y)= \begin{cases} F(x)^nF(y), & x\le y \\ F(y)^{n+1}, & x>y \end{cases}$

\end{document}
